% GENERATED FILE. Edit canonical lessons, then run npm run generate:books.
% canonical-source-hash: fnv1a64:da20dad45163b058
% canonical-lessons: JA-C316-monogatari, JA-C316-dentou, JA-C316-hotondo, JA-C316-kanjiru, JA-C316-motto, JA-R316-review-312-316, JA-R316-review-final-316

\chapter{A Story to Pass On}
\label{ch:ja-a-story-to-pass-on}
\hlchaptermodality{316}

\begin{chapteropening}
\textbf{By the end of this chapter:} \emph{I can say what old stories and customs make me feel.}

Complete the last lesson of chapter 316: I can say what old stories and customs make me feel.
\end{chapteropening}

\section[\texorpdfstring{\ja{ものがたり} (monogatari)}{ものがたり (monogatari)}]{\ja{ものがたり} (monogatari) — a story; a tale}

\label{lesson:JA-C316-monogatari}



\begin{quote}
Before the new one: say the Japanese for here and there, then the Japanese for to keep watch.
\end{quote}

\subsection*{You'll want to know: \ja{ものがたり}}
\textbf{\ja{ものがたり}} — \emph{monogatari} — \textquotedblleft{}a story; a tale\textquotedblright{}.

\begin{quote}
\textbf{\ja{これは} \ja{むかしの} \ja{ものがたりです。}} — \emph{kore wa mukashi no monogatari desu} — \textquotedblleft{}This is a story from long ago.\textquotedblright{}
\end{quote}

\begin{cousinweb}
\textbf{\ja{ものがたり}} (\emph{monogatari}) is a native Japanese word (wago), inherited rather than borrowed from Chinese, built from \emph{mono}, \textquotedblleft{}thing\textquotedblright{}, and \emph{gatari}, \textquotedblleft{}telling\textquotedblright{} (from \emph{kataru}) (JMdict, KANJIDIC2).
\end{cousinweb}

\begin{grammarlens}[title={What you've built}]
One more word for telling what happened.
\end{grammarlens}

\subsection*{Your turn}
\begin{itemize}
\raggedright
  \item \emph{Say it:} \emph{monogatari}
  \item \emph{Say it:} \emph{monogatari}, once more
  \item \emph{Say it:} \emph{kore wa mukashi no monogatari desu} — This is a story from long ago.
  \item \emph{Recall:} say the Japanese for here and there, then the Japanese for to keep watch, then say \emph{monogatari} again
\end{itemize}

\begin{tcolorbox}[breakable,colback=teal!4,colframe=teal!35!black,title={Before you move on}]
What does \ja{ものがたり} mean? (\textquotedblleft{}A story; a tale\textquotedblright{}.) Say it once more.
\end{tcolorbox}
\section[\texorpdfstring{\ja{でんとう} (dentō)}{でんとう (dentō)}]{\ja{でんとう} (dentō) — a tradition}

\label{lesson:JA-C316-dentou}



\begin{quote}
Before the new one: say the Japanese for to keep watch, then the Japanese for a story.
\end{quote}

\subsection*{You'll want to know: \ja{でんとう}}
\textbf{\ja{でんとう}} — \emph{dentō} — \textquotedblleft{}a tradition\textquotedblright{}.

\begin{quote}
\textbf{\ja{この} \ja{まちには} \ja{ふるい} \ja{でんとうが} \ja{あります。}} — \emph{kono machi ni wa furui dentō ga arimasu} — \textquotedblleft{}This town has old traditions.\textquotedblright{}
\end{quote}

\begin{cousinweb}
\textbf{\ja{でんとう}} (\emph{dentō}) is Sino-Japanese (kango): two Chinese characters read with their on'yomi, readings borrowed from Chinese. They are \emph{den}, \textquotedblleft{}transmit\textquotedblright{} (Middle Chinese \emph{drjwen}) and \emph{tō}, \textquotedblleft{}relationship\textquotedblright{} (\emph{thowngH}) (KANJIDIC2; Middle Chinese in Baxter's transcription).
\end{cousinweb}

\begin{grammarlens}[title={What you've built}]
One more word for telling what happened.
\end{grammarlens}

\subsection*{Your turn}
\begin{itemize}
\raggedright
  \item \emph{Say it:} \emph{dentō}
  \item \emph{Say it:} \emph{dentō}, once more
  \item \emph{Say it:} \emph{kono machi ni wa furui dentō ga arimasu} — This town has old traditions.
  \item \emph{Recall:} say the Japanese for to keep watch, then the Japanese for a story, then say \emph{dentō} again
\end{itemize}

\begin{tcolorbox}[breakable,colback=teal!4,colframe=teal!35!black,title={Before you move on}]
What does \ja{でんとう} mean? (\textquotedblleft{}A tradition\textquotedblright{}.) Say it once more.
\end{tcolorbox}
\section[\texorpdfstring{\ja{ほとんど} (hotondo)}{ほとんど (hotondo)}]{\ja{ほとんど} (hotondo) — almost; nearly all}

\label{lesson:JA-C316-hotondo}



\begin{quote}
Before the new one: say the Japanese for a story, then the Japanese for a tradition.
\end{quote}

\subsection*{You'll want to know: \ja{ほとんど}}
\textbf{\ja{ほとんど}} — \emph{hotondo} — \textquotedblleft{}almost; nearly all\textquotedblright{}.

\begin{quote}
\textbf{\ja{ともだちの} \ja{なまえは} \ja{ほとんど} \ja{わすれました。}} — \emph{tomodachi no namae wa hotondo wasuremashita} — \textquotedblleft{}I have forgotten nearly all my friends' names.\textquotedblright{}
\end{quote}

\begin{cousinweb}
\textbf{\ja{ほとんど}} (\emph{hotondo}) is a native Japanese word (wago), inherited rather than borrowed from Chinese, and it is usually written in kana (JMdict).
\end{cousinweb}

\begin{grammarlens}[title={What you've built}]
One more word for telling what happened.
\end{grammarlens}

\subsection*{Your turn}
\begin{itemize}
\raggedright
  \item \emph{Say it:} \emph{hotondo}
  \item \emph{Say it:} \emph{hotondo}, once more
  \item \emph{Say it:} \emph{tomodachi no namae wa hotondo wasuremashita} — I have forgotten nearly all my friends' names.
  \item \emph{Recall:} say the Japanese for a story, then the Japanese for a tradition, then say \emph{hotondo} again
\end{itemize}

\begin{tcolorbox}[breakable,colback=teal!4,colframe=teal!35!black,title={Before you move on}]
What does \ja{ほとんど} mean? (\textquotedblleft{}Almost; nearly all\textquotedblright{}.) Say it once more.
\end{tcolorbox}
\section[\texorpdfstring{\ja{かんじる} (kanjiru)}{かんじる (kanjiru)}]{\ja{かんじる} (kanjiru) — to feel; to sense}

\label{lesson:JA-C316-kanjiru}



\begin{quote}
Before the new one: say the Japanese for a tradition, then the Japanese for almost.
\end{quote}

\subsection*{You'll want to know: \ja{かんじる}}
\textbf{\ja{かんじる}} — \emph{kanjiru} — \textquotedblleft{}to feel; to sense\textquotedblright{}.

\begin{quote}
\textbf{\ja{いなかで} \ja{あきを} \ja{かんじました。}} — \emph{inaka de aki o kanjimashita} — \textquotedblleft{}In the countryside I could feel the autumn.\textquotedblright{}
\end{quote}

Polite \textbf{\ja{かんじます}}, \emph{kanjimasu}; an ichidan verb. Past (polite): \textbf{\ja{ははが} \ja{あきを} \ja{かんじました。}}, \emph{haha ga aki o kanjimashita}.

\begin{cousinweb}
\textbf{\ja{かんじる}} (\emph{kanjiru}) is a mixed word, native and Sino-Japanese: \emph{kan}, \textquotedblleft{}feeling\textquotedblright{}, is a Chinese character read with its on'yomi, a reading borrowed from Chinese (Middle Chinese \emph{komX}), and \emph{jiru}, \textquotedblleft{}a native verb ending\textquotedblright{}, is native Japanese (JMdict, KANJIDIC2; Middle Chinese in Baxter's transcription).
\end{cousinweb}

\begin{grammarlens}[title={What you've built}]
One more word for telling what happened.
\end{grammarlens}

\subsection*{Your turn}
\begin{itemize}
\raggedright
  \item \emph{Say it:} \emph{kanjiru}
  \item \emph{Say it:} \emph{kanjiru}, once more
  \item \emph{Say it:} \emph{inaka de aki o kanjimashita} — In the countryside I could feel the autumn.
  \item \emph{Recall:} say the Japanese for a tradition, then the Japanese for almost, then say \emph{kanjiru} again
\end{itemize}

\begin{tcolorbox}[breakable,colback=teal!4,colframe=teal!35!black,title={Before you move on}]
What does \ja{かんじる} mean? (\textquotedblleft{}To feel; to sense\textquotedblright{}.) Say it once more.
\end{tcolorbox}
\section[\texorpdfstring{\ja{もっと} (motto)}{もっと (motto)}]{\ja{もっと} (motto) — more; even more}

\label{lesson:JA-C316-motto}



\begin{quote}
Before the new one: say the Japanese for almost, then the Japanese for to feel.
\end{quote}

\subsection*{You'll want to know: \ja{もっと}}
\textbf{\ja{もっと}} — \emph{motto} — \textquotedblleft{}more; even more\textquotedblright{}.

\begin{quote}
\textbf{\ja{もっと} \ja{むかしばなしを} \ja{ききたいです。}} — \emph{motto mukashibanashi o kikitai desu} — \textquotedblleft{}I want to hear more old stories.\textquotedblright{}
\end{quote}

\begin{cousinweb}
\textbf{\ja{もっと}} (\emph{motto}) is a native Japanese word (wago), inherited rather than borrowed from Chinese, and it is usually written in kana (JMdict).
\end{cousinweb}

\begin{grammarlens}[title={What you've built}]
One more word for telling what happened.
\end{grammarlens}

\subsection*{Your turn}
\begin{itemize}
\raggedright
  \item \emph{Say it:} \emph{motto}
  \item \emph{Say it:} \emph{motto}, once more
  \item \emph{Say it:} \emph{motto mukashibanashi o kikitai desu} — I want to hear more old stories.
  \item \emph{Recall:} say the Japanese for almost, then the Japanese for to feel, then say \emph{motto} again
\end{itemize}

\begin{tcolorbox}[breakable,colback=teal!4,colframe=teal!35!black,title={Before you move on}]
What does \ja{もっと} mean? (\textquotedblleft{}More; even more\textquotedblright{}.) Say it once more.
\end{tcolorbox}
\section[(dialogue)]{Looking back: from History to A Story to Pass On}

\label{lesson:JA-R316-review-312-316}



\begin{quote}
Say the Japanese for history and an era.
\end{quote}

\subsection*{Your turn}
Say these aloud:

\begin{itemize}
\raggedright
  \item history, an era, the past, formerly, at that time
  \item a record, to prosper, to perish, the other day, ruins
  \item completely, before one knows it, before long, to decrease, to think of
  \item thereupon, however, but, here and there, to keep watch
  \item a story, a tradition, almost, to feel, more
\end{itemize}

\begin{tcolorbox}[breakable,colback=teal!4,colframe=teal!35!black,title={Before you move on}]
History? (\textbf{\ja{れきし}}, \emph{rekishi}.) An era? (\textbf{\ja{じだい}}, \emph{jidai}.) The past? (\textbf{\ja{かこ}}, \emph{kako}.) Formerly? (\textbf{\ja{いぜん}}, \emph{izen}.) At that time? (\textbf{\ja{とうじ}}, \emph{tōji}.) A record? (\textbf{\ja{きろく}}, \emph{kiroku}.) To prosper? (\textbf{\ja{さかえる}}, \emph{sakaeru}.) To perish? (\textbf{\ja{ほろびる}}, \emph{horobiru}.) The other day? (\textbf{\ja{せんじつ}}, \emph{senjitsu}.) Ruins? (\textbf{\ja{いせき}}, \emph{iseki}.) Completely? (\textbf{\ja{すっかり}}, \emph{sukkari}.) Before one knows it? (\textbf{\ja{いつのまにか}}, \emph{itsunomanika}.) Before long? (\textbf{\ja{やがて}}, \emph{yagate}.) To decrease? (\textbf{\ja{へる}}, \emph{heru}.) To think of? (\textbf{\ja{おもいつく}}, \emph{omoitsuku}.) Thereupon? (\textbf{\ja{すると}}, \emph{suruto}.) However? (\textbf{\ja{ところが}}, \emph{tokoroga}.) But? (\textbf{\ja{しかし}}, \emph{shikashi}.) Here and there? (\textbf{\ja{ところどころ}}, \emph{tokorodokoro}.) To keep watch? (\textbf{\ja{みはる}}, \emph{miharu}.) A story? (\textbf{\ja{ものがたり}}, \emph{monogatari}.) A tradition? (\textbf{\ja{でんとう}}, \emph{dentō}.) Almost? (\textbf{\ja{ほとんど}}, \emph{hotondo}.) To feel? (\textbf{\ja{かんじる}}, \emph{kanjiru}.) More? (\textbf{\ja{もっと}}, \emph{motto}.)
\end{tcolorbox}
\section[(dialogue)]{A Story to Pass On — once more}

\label{lesson:JA-R316-review-final-316}



\begin{quote}
Say the Japanese for a story and a tradition.
\end{quote}

\subsection*{Your turn}
\begin{itemize}
\raggedright
  \item \emph{Say it:} a story, a tradition, almost, to feel, more
\end{itemize}

\begin{tcolorbox}[breakable,colback=teal!4,colframe=teal!35!black,title={Before you move on}]
A story? (\textbf{\ja{ものがたり}}, \emph{monogatari}.) A tradition? (\textbf{\ja{でんとう}}, \emph{dentō}.) Almost? (\textbf{\ja{ほとんど}}, \emph{hotondo}.) To feel? (\textbf{\ja{かんじる}}, \emph{kanjiru}.) More? (\textbf{\ja{もっと}}, \emph{motto}.)
\end{tcolorbox}
